\chapter{Tonsillitis}

\section*{Definition}
Tonsillitis is inflammation of the palatine tonsils, typically of infectious origin, diagnosed clinically by sore throat, red and swollen tonsils, with or without exudate. The Centor criteria (fever, tonsillar exudate, tender cervical lymphadenopathy, absence of cough) stratify likelihood of group A streptococcal infection.

\section*{What Happens in Your Body}
Infectious pathogens (most commonly viruses or group A \textit{Streptococcus}) breach the tonsillar crypt epithelium, triggering a local inflammatory cascade with neutrophil infiltration, widening of blood vessels, and swelling. The immune response within tonsillar lymphoid tissue produces exudate and systemic symptoms through immune signaling proteins release.

\section*{Causes}
\begin{itemize}
  \item \textbf{Viral:} Adenovirus, Epstein-Barr virus (EBV), influenza, parainfluenza, enterovirus, rhinovirus
  \item \textbf{Bacterial:} Group A \textit{Streptococcus pyogenes} (most common bacterial), \textit{Neisseria gonorrhoeae}, \textit{Corynebacterium diphtheriae}
  \item \textbf{Fungal:} \textit{Candida} species (immunocompromised hosts)
  \item \textbf{Other:} Smoking, environmental irritants, chronic mouth breathing
\end{itemize}

\section*{When to See a Doctor}
See your doctor if:
\begin{itemize}
  \item Severe throat pain with difficulty swallowing or drooling
  \item High fever (greater than 101.3\textdegree{}F / 38.5\textdegree{}C) lasting over 48 hours
  \item Unilateral tonsillar swelling with uvular deviation suggesting peritonsillar collection of pus
  \item Suspicion of airway compromise or high-pitched breathing sound
  \item History of rheumatic fever or recurrent tonsillitis (more than 6 episodes/year)
\end{itemize}

\section*{Related Symptoms}
\begin{itemize}
  \item Sore throat, painful swallowing, cervical lymphadenopathy, fever, halitosis, referred otalgia
\end{itemize}
\textbf{Important:} Viral tonsillitis typically presents with cough, rhinorrhea, and hoarseness, whereas bacterial tonsillitis more often presents with high fever and exudate.

\section*{Self-Care / Home Management}
\begin{itemize}
  \item Warm salt water gargles (1/2 teaspoon salt in 8 oz warm water) multiple times daily
  \item Increased fluid intake (warm tea with honey, cold beverages) for comfort and hydration
  \item Acetaminophen or ibuprofen for pain and fever reduction
  \item Soft, cool foods (ice cream, yogurt, smoothies) to minimize swallowing discomfort
  \item Humidifier to soothe throat mucosa and reduce irritation
\end{itemize}

\section*{How Your Doctor Will Evaluate You}
\begin{itemize}
  \item Modified Centor score (0--4): scores of 2 or higher warrant rapid antigen detection test (RADT) or throat culture
  \item Rapid strep test: high specificity (greater than 95\%) but variable sensitivity; confirm negative results in children with throat culture
  \item Complete blood count to differentiate viral from bacterial infection if diagnosis uncertain
  \item Heterophile antibody (Monospot) test if EBV infectious mononucleosis is suspected
\end{itemize}

\section*{Treatment Options}
\begin{itemize}
  \item Supportive care (analgesics, hydration, rest) is sufficient for viral tonsillitis
  \item Antibiotics for confirmed group A streptococcal infection: penicillin V or amoxicillin for 10 days
  \item Corticosteroids (single dose of dexamethasone 0.6 mg/kg) for severe pain or airway swelling
  \item Tonsillectomy considered for recurrent infections (more than 7 episodes/year, more than 5/year for 2 years, or more than 3/year for 3 years)
  \item Drainage and antibiotics for peritonsillar collection of pus; airway monitoring if compromised
\end{itemize}

\section*{References}
\begin{thebibliography}{99}
\bibitem{tonsillitis:ref1} Windfuhr JP, Toepfner N, Steffen G, et al. "Clinical practice guideline: Tonsillitis I. Diagnostics and non-surgical management." \textit{European Archives of Oto-Rhino-Laryngology}, 2016;273(4):973--987.
\bibitem{tonsillitis:ref2} Shulman ST, Bisno AL, Clegg HW, et al. "Clinical practice guideline for the diagnosis and management of group A streptococcal pharyngitis." \textit{Clinical Infectious Diseases}, 2012;55(10):e86--e102.
\bibitem{tonsillitis:ref3} Baugh RF, Archer SM, Mitchell RB, et al. "Clinical practice guideline: Tonsillectomy in children." \textit{Otolaryngology--Head and Neck Surgery}, 2011;144(1 Suppl):S1--S30.
\bibitem{tonsillitis:ref4} Parrott GL, Shaffer AD, Stapleton AL, et al. "Tonsillitis: A review of current management." \textit{Journal of the American Academy of Physician Assistants}, 2020;33(10):34--39.
\bibitem{tonsillitis:ref5} Burton MJ, Glasziou PP, Chong LY, et al. "Tonsillectomy versus non-surgical treatment for chronic/recurrent acute tonsillitis." \textit{Cochrane Database of Systematic Reviews}, 2014;11:CD001802.
\bibitem{tonsillitis:ref6} Del Mar CB, Glasziou PP, Spinks AB. "Antibiotics for sore throat." \textit{Cochrane Database of Systematic Reviews}, 2006;4:CD000023.
\end{thebibliography}
